\documentclass[11pt]{article}
\usepackage{amsmath}
\usepackage{amssymb}
\usepackage{amsthm}
\usepackage[margin=1in]{geometry}

\newtheorem{theorem}{Theorem}
\newtheorem{proposition}[theorem]{Proposition}
\newtheorem{corollary}[theorem]{Corollary}
\theoremstyle{definition}
\newtheorem{definition}[theorem]{Definition}
\newtheorem{remark}[theorem]{Remark}
\newtheorem{example}[theorem]{Example}

\newcommand{\Fix}{\operatorname{Fix}}
\newcommand{\id}{\mathrm{id}}

\title{A disproof of conjecture 00000002601:\\
the maximum number of fixed points of a monotone self-map\\
of a finite lattice is the lattice size, not the chain length}
\author{earthking11}
\date{2026-09-13}

\begin{document}
\maketitle

\begin{abstract}
Conjecture 00000002601 asserts that the maximum number of fixed points of a
monotone self-map of an $n$-element finite lattice equals the length of a
maximal chain, that this maximum is attained on distributive lattices, and
that on non-distributive lattices the maximum drops strictly, the drop being
the minimal number of embedded $N_5$ sublattices. We show the conjecture is
false. The identity map is monotone and fixes all $n$ elements, so the maximum
is exactly $n$ --- the trivial upper bound --- on every finite lattice. On the
four-element distributive lattice $B_2$ the maximum is $4$, while a maximal
chain $\bot < a < \top$ has $3$ elements (equivalently $2$ edges). The claimed
equality therefore fails already on a distributive lattice, and it fails under
both readings of ``chain length''. The strict drop on non-distributive
lattices is likewise false: on $N_5$ the identity still fixes all $5$
elements, so the drop is $0$. All claims are verified by exhaustive
enumeration and, independently, in core Lean~4.
\end{abstract}

\section{The conjecture and its refutation}

\begin{definition}[monotone self-map, fixed-point count]
Let $L$ be a finite lattice with order $\le$. A self-map $f \colon L \to L$ is
\emph{monotone} (order-preserving) if $x \le y$ implies $f(x) \le f(y)$ for all
$x,y \in L$. Its fixed-point set is $\Fix(f) = \{x \in L : f(x) = x\}$.
\end{definition}

The conjecture under consideration is the following.

\begin{quote}
\textbf{Conjecture 00000002601.} \emph{The maximum number of fixed points of a
monotone map on an $n$-element finite lattice equals the length of a maximal
chain, attained on distributive lattices; the maximum drops strictly on
non-distributive lattices, the drop being the minimal number of embedded $N_5$
sublattices.}
\end{quote}

The length of a chain admits two conventions: the number of its elements, or
the number of its edges (covering pairs). The conjecture is ambiguous between
them; we refute it under both. The main observation is utterly elementary.

\begin{theorem}[identity bound]
\label{thm:identity}
For every finite lattice $L$ with $|L| = n$, the identity map $\id_L$ is
monotone and satisfies $\Fix(\id_L) = L$, so $|\Fix(\id_L)| = n$. Since no
self-map of $L$ can fix more than $n$ elements, the maximum number of fixed
points over all monotone self-maps of $L$ is exactly $n$, the trivial upper
bound.
\end{theorem}

\begin{proof}
For $x \le y$ we have $\id_L(x) = x \le y = \id_L(y)$, so $\id_L$ is monotone.
Every element is fixed, hence $|\Fix(\id_L)| = n$. The number of fixed points
of any self-map is at most $|L| = n$. Therefore the maximum equals $n$.
\end{proof}

Consequently the quantity the conjecture tries to compute is trivially $n$ on
\emph{every} finite lattice, distributive or not. The conjectured equality
with the length of a maximal chain can only hold when a maximal chain has $n$
elements (equivalently $n-1$ edges), i.e.\ when the lattice is itself a chain.
As soon as $L$ is not a chain, some element lies off every maximal chain and
the equality fails. This already happens for the simplest non-chain
distributive lattice.

\begin{corollary}
\label{cor:general}
Let $L$ be a finite lattice of size $n$ that is not a chain. Then the maximum
number of fixed points of a monotone self-map of $L$ is $n$, which is strictly
larger than the number of elements of a maximal chain and strictly larger than
its number of edges. In particular Conjecture 00000002601 is false for every
finite lattice that is not a chain, and its ``attained on distributive
lattices'' rider does not rescue it.
\end{corollary}

\begin{proof}
By Theorem~\ref{thm:identity} the maximum is $n$. A maximal chain of $L$ has
at most $n-1$ elements, since $L$ is not a chain, and therefore at most $n-2$
edges; both are strictly less than $n$.
\end{proof}

\section{The explicit counterexample $B_2$}

Let $B_2 = \{0, a, b, 1\}$ be the Boolean lattice on two atoms, i.e.\ the
four-element distributive lattice with $0$ bottom, $1$ top, and $a, b$
incomparable. Its Hasse diagram is the square $0 < a < 1$, $0 < b < 1$.

\begin{example}
\label{ex:b2}
On $B_2$:
\begin{itemize}
\item the identity map is monotone and has $4$ fixed points;
\item a maximal chain is $0 < a < 1$, which has $3$ elements ($2$ edges);
\item hence the maximum number of fixed points ($4$) differs from the maximal
chain length under either convention: $4 \neq 3$ and $4 \neq 2$.
\end{itemize}
\end{example}

The seven non-identity elements of the lattice of all $36$ monotone self-maps
of $B_2$ fix at most $3$ elements, so the maximum $4$ is genuinely attained by
the identity and not by an accident of a degenerate map. In particular the
failure of the conjectured equality is not repaired by excluding the identity:
on $B_3$ (below) a non-identity monotone map still fixes $7$ elements, which
exceeds the maximal chain length $4$.

\section{Exhaustive computations}

We verified the relevant maxima by brute force. For completeness we record the
exact counts; all of them are reproduced by \texttt{reproduce.py} and, for
$B_2$, by the Lean development in \texttt{lean4/}.

\subsection{$B_2$: all $4^4 = 256$ self-maps}

Enumerating all $256$ self-maps of $B_2$ and retaining the monotone ones gives
exactly $36$ monotone maps. The maximum of $|\Fix(f)|$ over them is $4$,
attained by the identity; the maximum over non-identity monotone maps is $3$.
The maximal chain has $3$ elements and $2$ edges. Hence
\[
  4 = \max_f |\Fix(f)| \;\neq\; 3 = \#\{\text{elements of a maximal chain}\},
\]
and likewise $4 \neq 2$ for the edge count.

\subsection{$B_3$: all $8^8 = 16\,777\,216$ self-maps, enumerated by pruning}

Let $B_3$ be the Boolean lattice on three atoms, of size $8$. A naive scan of
all $8^8 = 16\,777\,216$ self-maps is possible in principle but wasteful; the
script instead enumerates the monotone maps directly by depth-first search,
assigning the image of each element in a topological order and pruning any
partial assignment that violates monotonicity. This enumerates exactly the
monotone self-maps and finds
\[
  \#\{\text{monotone self-maps of } B_3\} = 8000,
\]
with maximum fixed-point count $8$ (the identity) and maximum over
non-identity monotone maps $7$. A maximal chain of $B_3$ has $4$ elements
($3$ edges). Thus $8 \neq 4$ and $8 \neq 3$.

\begin{remark}[non-identity maps]
\label{rem:nonid}
On $B_3$, the monotone map that sends one coatom to the top element and fixes
everything else is not the identity and fixes $7$ elements. Since $7 > 4$ and
$7 > 3$, excluding the identity does not restore the conjecture.
\end{remark}

\section{The non-distributive lattice $N_5$}

Let $N_5$ be the five-element non-distributive lattice with bottom $0$, top
$1$, and middle elements $a, b, c$ satisfying
\[
  0 < a < b < 1, \qquad 0 < a < c < 1,
\]
with $b$ and $c$ incomparable. The identity is monotone and fixes all $5$
elements, so the maximum number of fixed points is $5$, while a maximal chain
$0 < a < b < 1$ has $4$ elements ($3$ edges). Consequently
\[
  \max_f |\Fix(f)| = 5 \neq 4 \quad\text{and}\quad 5 \neq 3 .
\]

More importantly for the second half of the conjecture: the maximum on $N_5$
is $5 = |N_5|$, exactly the same trivial value as on \emph{any} five-element
lattice, distributive or not. So the maximum does \emph{not} drop strictly on
the non-distributive lattice $N_5$; the drop is $0$. Any clause tying a strict
drop to non-distributivity, or to the number of embedded $N_5$ sublattices, is
therefore false as stated: the conjectured quantity is blind to
distributivity. (For reference, exhaustive enumeration shows $N_5$ has $133$
monotone self-maps, the maximum among non-identity ones being $4$; but the
identity alone already realizes the true maximum $5$.)

\section{The elements-versus-edges ambiguity}

The conjecture says ``the length of a maximal chain'' without specifying
whether a chain of $k$ elements has length $k$ or $k-1$. We stress that the
refutation is independent of this choice:
\begin{center}
\begin{tabular}{lcccc}
\hline
Lattice & $\max_f |\Fix(f)|$ & chain elements & chain edges & refuted? \\
\hline
$B_2$ (distr., $n=4$) & $4$ & $3$ & $2$ & yes (both) \\
$B_3$ (distr., $n=8$) & $8$ & $4$ & $3$ & yes (both) \\
$N_5$ (non-distr., $n=5$) & $5$ & $4$ & $3$ & yes (both) \\
\hline
\end{tabular}
\end{center}
In every case the maximum number of fixed points equals the lattice size $n$
(Theorem~\ref{thm:identity}) and is strictly larger than both candidate chain
lengths. The conjectured equality fails on the distributive lattices $B_2$ and
$B_3$ as well, so the ``attained on distributive lattices'' part is false too.

\section{Formal verification}

The file \texttt{lean4/Main.lean} formalizes the $B_2$ refutation in core Lean~4
(\texttt{import Std}, no Mathlib, no \texttt{sorry}, no \texttt{axiom}, no
\texttt{native\_decide}). It encodes $B_2$ as $\mathrm{Fin}\,4$, defines
monotonicity, proves that the identity is monotone and fixes all four elements,
and proves by exhaustive computation over the $256$ functions that
\[
  \texttt{maxFixedPoints} = 4, \qquad
  \texttt{maxChainLen} = 3, \qquad
  4 > 3,
\]
culminating in \texttt{conjecture\_00000002601\_false}. A companion
\texttt{Check.lean} prints the axioms of every theorem; none depends on
\texttt{sorryAx} or \texttt{ofReduceBool}. The script \texttt{reproduce.py}
independently reproduces all numerical claims above (including $B_3$ and
$N_5$) using only the Python standard library.

\begin{thebibliography}{9}
\bibitem{tarski}
A.~Tarski, \emph{A lattice-theoretical fixpoint theorem and its applications},
Pacific J. Math. \textbf{5} (1955), 285--309.
\bibitem{grätzer}
G.~Grätzer, \emph{Lattice Theory: Foundation}, Birkhäuser, 2011.
\end{thebibliography}

\end{document}
